%Version 3.1 December 2024
% See section 11 of the User Manual for version history
%
%%%%%%%%%%%%%%%%%%%%%%%%%%%%%%%%%%%%%%%%%%%%%%%%%%%%%%%%%%%%%%%%%%%%%%
%%                                                                 %%
%% Please do not use \input{...} to include other tex files.       %%
%% Submit your LaTeX manuscript as one .tex document.              %%
%%                                                                 %%
%% All additional figures and files should be attached             %%
%% separately and not embedded in the \TeX\ document itself.       %%
%%                                                                 %%
%%%%%%%%%%%%%%%%%%%%%%%%%%%%%%%%%%%%%%%%%%%%%%%%%%%%%%%%%%%%%%%%%%%%%

\documentclass[pdflatex,sn-mathphys-num]{sn-jnl}% Math and Physical Sciences Numbered Reference Style

%%%% Standard Packages
\usepackage{graphicx}%
\usepackage{multirow}%
\usepackage{amsmath,amssymb,amsfonts}%
\usepackage{amsthm}%
\usepackage{mathrsfs}%
\usepackage[title]{appendix}%
\usepackage{xcolor}%
\usepackage{textcomp}%
\usepackage{manyfoot}%
\usepackage{booktabs}%
\usepackage{algorithm}%
\usepackage{algorithmicx}%
\usepackage{algpseudocode}%
\usepackage{listings}%
%%%%

\theoremstyle{thmstyleone}%
\newtheorem{theorem}{Theorem}%
\newtheorem{proposition}[theorem]{Proposition}%

\theoremstyle{thmstyletwo}%
\newtheorem{example}{Example}%
\newtheorem{remark}{Remark}%

\theoremstyle{thmstylethree}%
\newtheorem{definition}{Definition}%

\raggedbottom

\begin{document}

\title[Autonomous Laundry Grasping]{Autonomous Laundry Grasping System combining Vision and Tactile Anomaly Detection}

\author*[1,2]{\fnm{First} \sur{Author}}\email{iauthor@gmail.com}
\author[2,3]{\fnm{Second} \sur{Author}}\email{iiauthor@gmail.com}
\equalcont{These authors contributed equally to this work.}
\author[1,2]{\fnm{Third} \sur{Author}}\email{iiiauthor@gmail.com}
\equalcont{These authors contributed equally to this work.}

\affil*[1]{\orgdiv{Department}, \orgname{Organization}, \orgaddress{\street{Street}, \city{City}, \postcode{100190}, \state{State}, \country{Country}}}

\abstract{Autonomous laundry handling remains a fundamental challenge in domestic robotics due to the extreme deformability of garments, significant intra-class shape variability, and the constrained nature of target environments like washing machines. While tactile sensing is crucial for grasp state detection, existing supervised approaches require exhaustive labeled failure data, which is intractable for highly deformable objects. Furthermore, vision-based grasping often suffers from inverse kinematics (IK) failures in confined spaces and false-positive detections in unstructured failure-recovery scenarios. In this paper, we propose a fully integrated autonomous laundry grasping system that tightly couples domain-conditioned visual perception with unsupervised tactile anomaly detection. For perception, we employ a lightweight, domain-conditioned segmentation network (TinyUNet) enhanced by an RGB gating mechanism to robustly localize grasp points while suppressing environmental false positives. To execute grasps in constrained spaces, we introduce a motion planning pipeline with orientation relaxation and a joint-based ``Blind Push'' fallback, effectively overcoming critical IK planning failures. During execution, grasp stability is continuously monitored using a Gated Recurrent Unit Autoencoder (GRU-AE) trained exclusively on normal tactile time-series data. Operating in a one-class paradigm, the tactile module detects grasp perturbations without requiring labeled failure modes, automatically triggering a closed-loop visual recovery pipeline upon failure. Extensive real-world experiments on a manipulator equipped with tactile sensors demonstrate that the proposed architecture significantly outperforms static baselines and end-to-end imitation learning models, achieving highly reliable laundry retrieval and deposition.}

\keywords{Deformable Object Manipulation, Tactile Anomaly Detection, Autonomous Laundry Grasping, Domestic Service Robotics, Grasp Failure Recovery}

\maketitle

% ============================================================
\section{Introduction}
\label{sec:intro}
% ============================================================

Autonomous laundry handling represents one of the most demanding challenges in domestic service robotics. Unlike rigid industrial objects, laundry items such as T-shirts, towels, and pants exhibit high intra-class shape variability, heterogeneous surface textures, and unpredictable deformation during grasping, making reliable manipulation significantly more difficult than in structured industrial settings~\cite{Billard2019, Hamajima1996}. A particularly critical sub-problem is grasp state detection: determining in real time whether the end-effector has established a stable hold on a garment, or whether a failure event such as slippage or insufficient contact has occurred.

Tactile sensing has emerged as a promising modality for grasp state detection, providing direct contact information that is largely invariant to the illumination variations and occlusion that limit vision-only approaches~\cite{Calandra2018, Yan2022ICRA}. However, existing tactile-based methods predominantly adopt a \textit{supervised classification} paradigm, requiring labeled examples of all grasp states — including failure modes — for training~\cite{YGC2026, Cui2020, Funabashi2022}. For rigid or quasi-rigid objects with well-defined failure modes, this is feasible. For highly deformable objects such as laundry, however, the space of possible contact patterns during both normal and failed grasps is vast and difficult to enumerate exhaustively, making the collection of comprehensive labeled failure datasets both costly and incomplete.

This limitation motivates an \textit{unsupervised one-class} approach, in which the model is trained exclusively on normal grasp data and anomalies are detected as deviations from the learned normal pattern. Such an approach requires no failure-mode labels, substantially reducing the data collection burden, and is inherently open to unseen failure modes that were not anticipated during training.

A further challenge is grasp point localization in the presence of clutter and deformability. Laundry items piled in a basket present no well-defined object boundaries, and the appearance of a fallen garment during failure recovery differs substantially from the same item in the basket, necessitating separate perceptual strategies for the two scenarios.

To address these challenges, we propose an integrated autonomous laundry grasping system that combines vision-based grasp point localization with tactile-based grasp anomaly detection. The vision module employs a lightweight domain-conditioned segmentation network that handles both initial grasping and failure recovery scenarios within a single model. To overcome False Positives in unstructured failure scenarios, we further integrate an RGB gating mechanism. Furthermore, a robust motion planning strategy incorporating ``Blind Push'' fallback behaviors ensures kinematic feasibility in constrained spaces. The tactile module employs a \textit{dual-channel} GRU-based Autoencoder (GRU-AE) trained in a one-class paradigm on normal grasp sequences. The primary channel (Ch-A) detects gradual grasp perturbations as temporal reconstruction anomalies in the 48-dimensional raw taxel time series, while the secondary channel (Ch-B) provides fast, frame-level response to abrupt contact changes via Mean Resultant Length (MRL) delta monitoring. Upon detecting a failure, the system automatically invokes the failure recovery pipeline and re-grasps the item.

The main contributions of this work are as follows:

\begin{enumerate}
    \item \textbf{Domain-conditioned segmentation with RGB Gating:}
    We propose a single lightweight segmentation network (TinyUNet) that handles two perceptually distinct scenarios via a task-conditioning channel. This is supplemented by an RGB Gating mechanism to suppress false positives on glossy surfaces, ensuring robust perception.

    \item \textbf{Dual-channel unsupervised tactile anomaly detection:}
    We propose a dual-channel GRU-AE framework trained exclusively on normal grasp data. Channel~A captures gradual slip through temporal reconstruction error on 48-dimensional raw taxel sequences, while Channel~B provides fast response to abrupt contact loss through frame-level MRL delta monitoring. The two channels operate with complementary latencies (530\,ms vs.\ 70\,ms) and are fused via an OR decision rule, achieving high recall at near-zero false positive rate.

    \item \textbf{Integrated autonomous grasping with Kinematic Fallback:}
    We present a complete system integrating vision and tactile modalities with a closed-loop failure recovery pipeline. To overcome Inverse Kinematics (IK) failures in constrained spaces, the system utilizes relaxed orientation tolerances and a joint-based ``Blind Push'' strategy, enabling reliable laundry retrieval from basket to washing machine.

    \item \textbf{Experimental validation:}
    We conduct comprehensive experiments comparing the proposed system against an end-to-end imitation learning baseline (ACT) at the full task level, and against multiple unsupervised anomaly detection baselines at the module level.
\end{enumerate}

The remainder of this paper is organized as follows. Section~\ref{sec:related} reviews related work. Section~\ref{sec:overview} describes the overall system architecture. Section~\ref{sec:vision} presents the vision-based localization and motion planning pipeline. Section~\ref{sec:anomaly} describes the tactile-based grasp anomaly detection module. Section~\ref{sec:experiments} reports experimental results, and Section~\ref{sec:conclusion} concludes the paper.

% ============================================================
\section{Related Work}
\label{sec:related}
% ============================================================
\subsection{Tactile-based Grasp State Detection}
Tactile sensing provides rich contact information that complements visual perception in grasp state detection tasks...

\subsection{Deformable and Fabric Object Grasping}
Grasping deformable objects such as fabric and clothing has received increasing attention in service robotics...

\subsection{Anomaly Detection for Robot Manipulation}
Anomaly detection methods have been applied to various robot manipulation tasks...

\subsection{End-to-end Imitation Learning for Grasping}
End-to-end imitation learning methods learn complete manipulation policies directly from human demonstrations...

% ============================================================
\section{System Overview}\label{sec:overview}
% ============================================================
This study presents an autonomous laundry grasping system that integrates vision-based grasp point localization with tactile-based grasp anomaly detection.

\subsection{Hardware Configuration}\label{sec:hardware}
The robotic platform consists of a \textbf{PIPER} robot arm with a parallel-jaw gripper. Each finger of the gripper is equipped with a \textbf{XELA uSPa44} tactile sensor. An RGB-D camera is mounted in an eye-in-hand configuration.

\subsection{Operation Sequence}\label{sec:sequence}
The grasping procedure proceeds as follows:
\begin{enumerate}
    \item \textbf{Initial localization:}
    The vision module runs TinyUNet in basket mode ($d_\text{dom} = 0.0$) to segment the laundry pile and compute a 3D grasp point.

    \item \textbf{Grasping:}
    The robot arm moves to the estimated grasp pose and closes the gripper around the target item. Tactile signals are acquired at 100\,Hz.

    \item \textbf{Grasp monitoring:}
    The dual-channel tactile anomaly detection module continuously evaluates the grasp state. If a stable grasp (SG) is confirmed, the arm proceeds to transport the item.

    \item \textbf{Failure recovery:}
    If a grasp perturbation (GP) is detected by either channel, the system invokes the failure recovery pipeline. TinyUNet is re-run in failure mode ($d_\text{dom} = 1.0$) to locate the fallen item, filtered by the RGB Gating module.

    \item \textbf{Target Deposition:}
    Once a stable grasp is verified and transport is complete, the robot localizes the opening of the washing machine. It navigates to a predefined mid-point approach pose and carefully inserts the item into the drum, continuously monitoring for unexpected collisions.
\end{enumerate}

% ============================================================
\section{Vision-based Grasp Point Localization and Execution}\label{sec:vision}
% ============================================================

Locating a graspable region on laundry items is challenging due to the highly deformable and texturally homogeneous nature of garments, and the significant appearance difference between two operational scenarios: items piled in a laundry basket, and items that have fallen to the floor following a failed grasp.

\subsection{Input Representation}
\label{sec:input_repr}
For each frame, we construct a 6-channel input tensor comprising Differential depth ($\Delta d$), Raw depth ($z$), RGB ($r, g, b$), and a Domain plane ($d_\text{dom} \in \{0.0,\, 1.0\}$).

\subsection{Network Architecture}
\label{sec:tinyunet}
We adopt a lightweight encoder-decoder network, \textbf{TinyUNet}, based on the U-Net architecture.

\subsection{Domain-conditioned Segmentation}
\label{sec:domain_cond}
By incorporating the domain plane channel $d_\text{dom}$, the network receives an explicit scalar signal indicating which segmentation target is expected at inference time.

\subsection{False Positive Rejection via RGB Gating}
\label{sec:rgb_gating}
While depth-based segmentation is effective for the structured basket scenario, the failure recovery scenario introduces significant environmental noise, such as glossy floor surfaces and complex lighting reflections. These artifacts often trigger false positive detections in the depth difference channel.

To mitigate this, we integrate an auxiliary RGB gating module based on a lightweight convolutional neural network (e.g., ResNet18). For each candidate region proposed by the TinyUNet, the RGB gating module evaluates the corresponding cropped color image to classify whether the region contains a true laundry item or empty background. The system accepts the grasp point only if the garment probability exceeds a predefined hysteresis threshold (e.g., $P(\text{cloth}) > 0.60$). This gating mechanism drastically reduces false grasp attempts on reflective floor surfaces during failure recovery.

\subsection{Grasp Point Computation and Execution Strategy}
\label{sec:grasp_point}
The 2D grasp point is computed as the centroid of the foreground mask and back-projected to a 3D point $\mathbf{p}^* \in \mathbb{R}^3$. Ideally, the gripper approaches the target vertically from above.

However, grasping in constrained environments—such as reaching deep into the corners of a laundry basket or navigating the narrow opening of a washing machine—frequently leads to Inverse Kinematics (IK) failures (e.g., MoveIt planning aborts) due to strict orientation constraints. To ensure kinematic feasibility, our motion planning pipeline actively relaxes the end-effector roll and pitch tolerances when a strict vertical approach fails.

Furthermore, if standard IK solutions remain unavailable even with relaxed orientation tolerances, the system automatically falls back to a joint-based ``Blind Push'' strategy. In this fallback mode, the robot relies on joint-space interpolation and physical compliance rather than precise Cartesian visual alignment. It pushes through minor obstructions while continuously monitoring the tactile sensors, allowing it to secure a grasp in highly constrained spaces where purely vision-guided Cartesian planning would consistently fail.

% ============================================================
\section{Tactile-based Grasp Anomaly Detection}
\label{sec:anomaly}
% ============================================================

\subsection{Grasp State Definition}
\label{sec:states}

We characterize the grasp interaction using three mutually exclusive states. During \textbf{No Contact (NC)}, the gripper is open or closing and the sensor reading reflects only inter-finger contact without an object. During \textbf{Stable Grasp (SG)}, the object is firmly held and the tactile signal is stationary or slowly varying. A \textbf{Grasp Perturbation (GP)} occurs when the object slips, is partially released, or experiences an abrupt external disturbance, resulting in a significant deviation of the tactile signal from the stable-grasp pattern.

The detection objective is to distinguish GP from SG in real time, using only models trained on SG sequences (one-class setting).

% ─────────────────────────────────────────────
\subsection{Tactile Feature Extraction}
\label{sec:features}
% ─────────────────────────────────────────────

The XELA uSPa44 sensor comprises a $4 \times 4$ array of taxels, each reporting three-axis force components $(f_x, f_y, f_z)$. At each timestep $t$, the raw taxel output of one finger is a 48-dimensional vector:
\begin{equation}
    \mathbf{r}_t = [f_x^{(1)}, f_y^{(1)}, f_z^{(1)},\; \ldots,\; f_x^{(16)}, f_y^{(16)}, f_z^{(16)}] \in \mathbb{R}^{48}.
    \label{eq:raw}
\end{equation}

\noindent\textbf{Mean Resultant Length (MRL).}
For Channel~B and interpretability analysis, we compute the scalar MRL as the mean of per-taxel 3-axis force magnitudes:
\begin{equation}
    \mathrm{MRL}_t = \frac{1}{16} \sum_{i=1}^{16} \left\| \mathbf{f}^{(i)}_t \right\|_2,
    \label{eq:mrl}
\end{equation}
where $\mathbf{f}^{(i)}_t = (f_x^{(i)}, f_y^{(i)}, f_z^{(i)})$ is the force vector of the $i$-th taxel at time $t$.

\noindent\textbf{Mean Z-channel (mean$_Z$).}
For heuristic and window-based baselines, we also use the mean of the 16 Z-axis readings:
\begin{equation}
    \overline{z}_t = \frac{1}{16} \sum_{i=1}^{16} f_z^{(i)},
    \label{eq:meanz}
\end{equation}
which carries a discriminative directional sign: during pre-contact (NC), the fingers press against each other and $\overline{z}_t < 0$; once an object is grasped, the object pushes back on the taxels and $\overline{z}_t > 0$.

All features are $z$-score normalized using statistics computed from the stable-grasp training sequences, excluding the first 50 frames (pre-contact warm-up).

% ─────────────────────────────────────────────
\subsection{Proposed: Dual-Channel GRU-AE}
\label{sec:dual_channel}
% ─────────────────────────────────────────────

Grasp perturbations manifest in two distinct temporal patterns: \emph{gradual slip}, characterized by a slow drift of the reconstructed tactile pattern, and \emph{abrupt drop}, caused by sudden loss of contact. A single detector optimized for one regime inevitably under-performs on the other. We therefore propose a dual-channel architecture in which the two channels operate in parallel with complementary response times and are fused via an OR decision rule.

\subsubsection{Channel A: Temporal Reconstruction (GRU-AE)}
\label{sec:channel_a}

Channel A captures gradual perturbations through reconstruction anomaly detection on sliding windows of the raw 48-dimensional taxel signal.

\paragraph{Architecture.}
The model is a GRU-based Autoencoder (GRU-AE) with a sequence encoder and sequence decoder, trained exclusively on stable-grasp sequences in a one-class setting.

The \textbf{Encoder} is a 2-layer GRU network with hidden dimension $H = 64$ that processes an input window $\mathbf{X} \in \mathbb{R}^{T \times 48}$ of $T = 50$ consecutive frames (500\,ms at 100\,Hz):
\begin{equation}
    \mathbf{h}_T = \mathrm{GRU}_{\mathrm{enc}}(\mathbf{X}) \in \mathbb{R}^{H},
    \label{eq:encoder}
\end{equation}
where $\mathbf{h}_T$ is the final hidden state of the last GRU layer.

The \textbf{Decoder} receives $\mathbf{h}_T$ as a fixed context vector, tiling it $T$ times as the decoder input, and reconstructs the input sequence through a symmetric 2-layer GRU followed by a linear projection:
\begin{equation}
    \hat{\mathbf{X}} = \mathrm{FC}\!\left(\mathrm{GRU}_{\mathrm{dec}}\!\left(\mathbf{h}_T \cdot \mathbf{1}_{T}^\top\right)\right) \in \mathbb{R}^{T \times 48}.
    \label{eq:decoder}
\end{equation}

\paragraph{Anomaly Score.}
The anomaly score for a window is the mean squared reconstruction error (MSE) averaged over all frames and all feature dimensions:
\begin{equation}
    s^A_t = \frac{1}{T \cdot D} \left\| \mathbf{X}_t - \hat{\mathbf{X}}_t \right\|_F^2,
    \label{eq:score_a}
\end{equation}
where $D = 48$ and $\|\cdot\|_F$ denotes the Frobenius norm.

\paragraph{Threshold.}
The detection threshold $\tau_A$ is set as the $p_A$-th percentile of $s^A$ values computed over all training windows from the stable-grasp dataset:
\begin{equation}
    \tau_A = \mathrm{Percentile}_{p_A}\!\left(\{s^A_t\}_{\text{train}}\right),
    \label{eq:threshold_a}
\end{equation}
with $p_A = 99$ for grasp monitoring and $p_A = 99.7$ for transport monitoring (where baseline motion introduces higher normal variability).

\paragraph{Detection Rule.}
A grasp perturbation is flagged by Channel~A when $n_A = 3$ consecutive windows exceed $\tau_A$. Using a sliding window with stride 1 (real-time inference), the minimum detection latency of Channel~A is:
\begin{equation}
    \tau^{\mathrm{lat}}_A = T + n_A - 1 = 52\ \text{frames} = 520\,\mathrm{ms}.
    \label{eq:latency_a}
\end{equation}

\subsubsection{Channel B: Frame-level MRL Delta}
\label{sec:channel_b}

Channel B is designed to detect abrupt, impulsive changes in contact force that manifest within a few frames — a regime in which the sliding-window approach of Channel~A is inherently slow due to its 50-frame buffer.

\paragraph{Feature and Reference.}
At each frame $t$, the scalar MRL $(\mathrm{MRL}_t)$ is computed as in Eq.~\eqref{eq:mrl}. A reference level is estimated as the rolling mean over the preceding $w = 3$ frames:
\begin{equation}
    r_t = \frac{1}{w}\sum_{k=t-w}^{t-1} \mathrm{MRL}_k.
    \label{eq:reference_b}
\end{equation}
The frame-level anomaly score is the absolute deviation from this reference:
\begin{equation}
    \delta_t = \left|\mathrm{MRL}_t - r_t\right|.
    \label{eq:delta_b}
\end{equation}

\paragraph{Threshold.}
The threshold $\tau_B$ is calibrated from normal grasp sequences as the $p_B$-th percentile of all $\delta_t$ values:
\begin{equation}
    \tau_B = \mathrm{Percentile}_{p_B}\!\left(\{\delta_t\}_{\text{train}}\right),
    \label{eq:threshold_b}
\end{equation}
with $p_B = 99$ for both grasp and transport phases.

\paragraph{Detection Rule.}
A grasp perturbation is flagged by Channel~B when $n_B = 4$ consecutive frames satisfy $\delta_t > \tau_B$. The minimum detection latency is:
\begin{equation}
    \tau^{\mathrm{lat}}_B = w + n_B - 1 = 6\ \text{frames} = 60\,\mathrm{ms}.
    \label{eq:latency_b}
\end{equation}

\subsubsection{Dual-Channel Decision Fusion}
\label{sec:fusion}

The two channels are fused using a logical OR rule: a grasp perturbation (GP) is declared if either channel independently triggers its detection criterion. Let $\mathcal{A}_t$ be the event that $n_A$ consecutive windows ending at time $t$ exceed $\tau_A$, and $\mathcal{B}_t$ be the event that $n_B$ consecutive frames ending at time $t$ exceed $\tau_B$. Then:
\begin{equation}
    \mathrm{GP}_t = \mathcal{A}_t \;\vee\; \mathcal{B}_t.
    \label{eq:fusion}
\end{equation}

Channel~A covers the slow, gradual slip regime with high precision, while Channel~B covers the fast, impulsive drop regime that Channel~A would detect too late. The parameters are summarized in Table~\ref{tab:dual_params}.

\begin{table}[h]
\caption{Dual-channel detector parameters}
\label{tab:dual_params}
\centering
\begin{tabular}{llcc}
\hline
Parameter & Description & Ch-A (GRU-AE) & Ch-B (MRL delta) \\
\hline
Input & Feature vector & 48D raw taxel & Scalar MRL \\
Temporal scope & Window / smoothing & $T=50$ frames & $w=3$ frames \\
Consec.\ threshold & $n_A$ / $n_B$ & 3 windows & 4 frames \\
Score threshold pct & $p_A$ / $p_B$ & 99 / 99.7 & 99 \\
Min.\ detection latency & $\tau^{\mathrm{lat}}$ & 520\,ms & 60\,ms \\
Target perturbation & Regime & Gradual slip & Abrupt drop \\
\hline
\end{tabular}
\end{table}

\subsection{Baseline Methods}
\label{sec:baselines}

We compare the proposed dual-channel GRU-AE against four baselines, all trained in the one-class setting on stable-grasp data only.

\subsubsection{Heuristic (mean$_Z$ threshold)}
\label{sec:heuristic}

The simplest baseline exploits the sign property of the Z-axis readings. As described in Section~\ref{sec:features}, $\overline{z}_t < 0$ during pre-contact and $\overline{z}_t > 0$ during stable grasping. A frame is declared anomalous if:
\begin{equation}
    \overline{z}_t < \zeta, \quad \zeta = 0.0,
    \label{eq:heuristic}
\end{equation}
i.e., the mean Z-channel drops back below zero after contact was established. The first 50 frames (pre-contact warm-up) are excluded from detection to avoid false alarms during gripper closure. No training is required.

\subsubsection{K-means and Fuzzy C-Means (FCM)}
\label{sec:kmeans_fcm}

Both methods cluster 2-dimensional window feature vectors extracted from stable-grasp sequences. For each sliding window of length $T_w = 50$ (stride 5), the feature vector is:
\begin{equation}
    \mathbf{v} = \left[\;\overline{\overline{z}}_w,\; \Delta\overline{z}_w\;\right] \in \mathbb{R}^2,
    \label{eq:win_feat}
\end{equation}
where $\overline{\overline{z}}_w$ is the window mean of $\overline{z}_t$ and $\Delta\overline{z}_w = \overline{z}^{(\text{2nd half})} - \overline{z}^{(\text{1st half})}$ is the intra-window delta, capturing both the level and the trend of the mean Z-channel.

Both methods fit $K = 3$ cluster centers $\{\mathbf{c}_k\}_{k=1}^3$ from normal training windows. The anomaly score for a test window is the minimum Euclidean distance to any cluster center:
\begin{equation}
    s^{\mathrm{cl}} = \min_{k=1,\ldots,3} \left\|\mathbf{v} - \mathbf{c}_k\right\|_2.
    \label{eq:cluster_score}
\end{equation}

Note that for FCM we use this distance-based score rather than the conventional membership-based score ($1 - \max_k u_k$), as the latter is bounded by the FCM geometry ($\approx 0.67$ for $K=3$) and therefore provides insufficient separation between normal and anomalous windows. The detection threshold is set at the 99th percentile of normal training window scores.

\subsubsection{Isolation Forest}
\label{sec:iforest}

Isolation Forest~\cite{Liu2008} is applied at the frame level, treating each frame independently without temporal modeling. The input feature is the 1-dimensional mean Z-channel scalar $\overline{z}_t$ (Eq.~\eqref{eq:meanz}). The model is fitted with 200 trees and contamination ratio 0.05 on stable-grasp frames (first 50 warm-up frames excluded). A frame is flagged as anomalous if the Isolation Forest assigns it a score of $-1$. A sequence is declared a GP if the binary anomaly flag is raised for at least $n$ consecutive frames ($n = 3$ in evaluation).

Note that using the 1D mean$_Z$ feature substantially outperforms using the full 16D Z-channel matrix, because the scalar mean captures the discriminative sign reversal ($\overline{z} < 0$) as a compact signal, whereas the 16D feature causes the forest to learn the heterogeneous normal distribution of individual Z channels, reducing sensitivity to anomalies.

% ============================================================
\section{Experiments}
\label{sec:experiments}
% ============================================================
\subsection{Experimental Setup}
All experiments were conducted using the PIPER robot arm equipped with a parallel-jaw gripper and XELA uSPa44 tactile sensors. Experiments were conducted with three categories of laundry items: T-shirts, towels, and pants.

We compare against \textbf{ACT (Action Chunking with Transformers)}, an end-to-end imitation learning method trained on 50 teleoperated demonstrations of the full laundry grasping and deposition task. While ACT learns a complete manipulation policy directly from demonstrations, it lacks an explicit failure detection or recovery mechanism, serving as a baseline to demonstrate the necessity of explicit anomaly detection and kinematic fallbacks.

\subsection{Tactile Anomaly Detection Evaluation}
\label{sec:tactile_eval}

\paragraph{Dataset.}
The grasp-phase dataset contains 51 stable-grasp sequences and 15 grasp-perturbation sequences, all recorded at 100\,Hz with the XELA uSPa44 sensor. The transport-phase dataset contains 43 stable-transport sequences and 12 transport-perturbation sequences. For all one-class methods, models are trained on stable-grasp (or stable-transport) sequences only; perturbation sequences are held out for evaluation.

\paragraph{Metrics.}
Each sequence is assigned a binary prediction (GP or SG) based on the detection criterion (consecutive anomaly frames/windows exceeding the respective threshold). We report Precision, Recall (TPR), F1, and AUROC. AUROC is approximated by varying the consecutive-window count $n \in \{1, 2, 3, 5\}$ and computing the trapezoidal area under the resulting TPR–FPR operating points.

\begin{table}[h]
\caption{Tactile anomaly detection performance — Grasp phase (51 SG / 15 GP sequences)}
\label{tab:tactile_eval_grasp}
\centering
\begin{tabular}{lcccc}
\hline
Method & Precision (\%) & Recall (\%) & F1 (\%) & AUROC \\
\hline
Heuristic (mean$_Z$)    & 51.7 & \textbf{100.0} & 68.1 & 0.86 \\
K-means                 & 56.0 & 93.3           & 69.6 & 0.88 \\
FCM                     & 66.7 & 93.3           & 77.8 & 0.93 \\
Isolation Forest        & 82.4 & 93.3           & 87.5 & 0.94 \\
GRU-AE Ch-A (ours)      & 93.8 & \textbf{100.0} & 96.8 & \textbf{0.99} \\
\textbf{GRU-AE Ch-A+B (ours)} & \textbf{100.0} & \textbf{100.0} & \textbf{100.0} & \textbf{0.99} \\
\hline
\end{tabular}
\end{table}

\begin{table}[h]
\caption{Tactile anomaly detection performance — Transport phase (43 SG / 12 GP sequences)}
\label{tab:tactile_eval_transport}
\centering
\begin{tabular}{lcccc}
\hline
Method & Precision (\%) & Recall (\%) & F1 (\%) & AUROC \\
\hline
Heuristic (mean$_Z$)    & [TODO] & [TODO] & [TODO] & [TODO] \\
K-means                 & [TODO] & [TODO] & [TODO] & [TODO] \\
FCM                     & [TODO] & [TODO] & [TODO] & [TODO] \\
Isolation Forest        & [TODO] & [TODO] & [TODO] & [TODO] \\
GRU-AE Ch-A (ours)      & [TODO] & [TODO] & [TODO] & [TODO] \\
\textbf{GRU-AE Ch-A+B (ours)} & [TODO] & [TODO] & [TODO] & [TODO] \\
\hline
\end{tabular}
\end{table}

The results in Table~\ref{tab:tactile_eval_grasp} demonstrate that GRU-AE outperforms all non-temporal baselines across all metrics, confirming that explicit temporal modeling of tactile signal evolution is essential for reliable anomaly detection during deformable object grasping. The addition of Channel~B (MRL delta) eliminates the remaining 2 false negatives from Ch-A alone at zero false positive cost, achieving perfect F1 on the grasp dataset. Among baselines, FCM and K-means benefit from the distance-based anomaly score (Eq.~\eqref{eq:cluster_score}), which provides an unbounded separation between normal and anomalous windows compared to the bounded membership-based score used in standard FCM.

\subsection{Vision Pipeline Evaluation}
\begin{table}[h]
\caption{TinyUNet segmentation performance per operational mode}
\label{tab:vision_eval}
\centering
\begin{tabular}{lcc}
\hline
Mode & IoU (\%) & Dice (\%) \\
\hline
Basket ($d_\text{dom} = 0.0$)          & [TODO] & [TODO] \\
Failure recovery ($d_\text{dom} = 1.0$) & [TODO] & [TODO] \\
\hline
\end{tabular}
\end{table}

\subsection{Full Task Evaluation}
We evaluate the complete laundry grasping pipeline over 50 trials. Each trial begins with a laundry item placed in the basket and is considered successful if the item is securely deposited into the washing machine.

\begin{table}[h]
\caption{Full task success rate comparison (50 trials per method)}
\label{tab:task_eval}
\centering
\begin{tabular}{lccccc}
\hline
Method & T-shirt & Towel & Pants & Total TSR (\%) & RSR (\%) \\
\hline
ACT (50 demos)
  & [?]/20 & [?]/15 & [?]/15 & $\sim$70 & — \\
\textbf{Ours (first attempt)}
  & [?]/20 & [?]/15 & [?]/15 & [TODO] & — \\
\textbf{Ours (with recovery)}
  & [?]/20 & [?]/15 & [?]/15 & [TODO] & [TODO] \\
\hline
\end{tabular}
\end{table}

As shown in Table~\ref{tab:task_eval}, the proposed system achieves a substantially higher task success rate than ACT across all laundry categories. The combination of depth-based localization, RGB gating, and flexible kinematic fallbacks allows the system to consistently secure grasps even in highly constrained states, and further improves overall task success through the tactile failure recovery pipeline.

\section{Conclusion}
\label{sec:conclusion}
We presented an autonomous laundry grasping system integrating a domain-conditioned vision module, RGB gating, and a dual-channel one-class tactile anomaly detector. Channel~A (GRU-AE) captures gradual slip through temporal reconstruction error on 48-dimensional raw taxel sequences, while Channel~B (MRL delta) provides fast, frame-level response to abrupt contact loss at 60\,ms latency. The OR-fusion of the two channels achieved perfect detection at zero false positive rate on the grasp-phase dataset, outperforming all non-temporal baselines. The inclusion of flexible motion planning strategies, such as ``Blind Push'' and orientation relaxation, enabled robust operation in highly constrained spaces like washing machine drums. Experimental results confirmed the superior performance of the proposed architecture over both static baselines and end-to-end learning approaches.

\bibliography{isr_references}
\end{document}
